% Avoid leaving only a heading and two lines at the foot of the preceding column.
\Needspace{14\baselineskip}
\section{Graph models and core observables}
\label{sec:models-and-observables}

\subsection{The 3-core and its size}
\label{sec:core-and-observables}

To find the 3-core, remove every vertex with fewer than three remaining
neighbours, then repeat: vertices that initially had enough neighbours may
lose them during the deletions. A simple undirected graph $G=(V,E)$ has
vertex set $V$ and edge set $E$; edges join distinct vertices and each pair
has at most one edge.
Write $n=|V|\geq1$ for its initial number of vertices. For a subset $S\subseteq V$,
the induced subgraph $G[S]$ keeps the vertices in $S$ and every edge of $G$
between them. Let $d_{G[S]}(v)$ count the neighbours of $v$ in this subgraph.
The simultaneous deletion rounds are
\begin{equation}
  S_0=V,\qquad
  S_{t+1}=\{v\in S_t:d_{G[S_t]}(v)\geq3\},
  \label{eq:core-pruning}
\end{equation}
where $t=0,1,\ldots$ counts rounds. The sets decrease and eventually stop
changing. If $S_\infty$ is the final set, then
$C_3(G)=G[S_\infty]$ is the 3-core~\cite{janson2007kcore}.

This procedure gives the unique largest induced subgraph with minimum degree
at least three. Indeed, if $G[W]$ already has that minimum degree, none of its
vertices can be the first vertex of $W$ removed: its three neighbours in $W$
would still be present. Thus every such $W$ remains inside the final set,
which itself meets the degree condition. This also proves that deleting
eligible vertices one at a time, in any order, gives the same answer.
The core may be disconnected; we count all its vertices, rather than only
its largest connected component.

The implementation uses a queue to avoid rescanning the whole graph after
each deletion. Initially queue all vertices of degree below three.
Remove the first queued vertex and its incident edges, update the degrees
of its remaining neighbours, and queue a neighbour when its degree falls
from three to two. Stop when the queue is empty. Each vertex is queued at
most once, giving work proportional to $|V|+|E|$.
The input graph is left unchanged. \Cref{fig:core-pruning} illustrates why a
single pass would be insufficient. For the configuration multigraph defined
in \cref{sec:configuration-model}, the same rule uses incident edge ends
instead of distinct neighbours, with the degree convention specified there.

\begin{figure}[b]
  \centering
  % Deterministic pruning example. The solid teal K_4 survives; dashed edges
% touch vertices deleted in the first two simultaneous pruning rounds.
% This source is a single-column graphic; its caption belongs in the section.
\begingroup
\definecolor{coreteal}{RGB}{33,112,117}
\definecolor{prunegray}{RGB}{150,156,160}
\definecolor{pruneorange}{RGB}{182,112,49}
\begin{tikzpicture}[
  x=1.2cm, y=1.05cm,
  vertex/.style={circle, minimum size=5.5mm, inner sep=0pt, font=\small},
  retained/.style={vertex, draw=coreteal, fill=coreteal!10, text=coreteal,
    line width=0.8pt},
  first/.style={vertex, draw=prunegray, fill=prunegray!12, text=black!75},
  second/.style={vertex, draw=pruneorange, fill=pruneorange!12,
    text=pruneorange},
  core edge/.style={draw=coreteal, line width=0.85pt},
  deleted edge/.style={draw=prunegray!85, dashed, line width=0.65pt}
]
  % Six edges on a,b,c,d form K_4, so every retained vertex has degree three.
  \coordinate (a) at (0,1);
  \coordinate (b) at (0,-0.5);
  \coordinate (c) at (1.5,1);
  \coordinate (d) at (1.5,-0.5);
  \coordinate (x) at (3,1);
  \coordinate (y) at (4.4,0.4);
  \coordinate (z) at (3.5,-0.7);
  \draw[core edge] (a)--(b)--(d)--(c)--(a);
  \draw[core edge] (a)--(d) (b)--(c);
  % Initially y,z have degree two, whereas x has degree three. Removing
  % y,z leaves x with only c,d, so x is removed in the next round.
  \draw[deleted edge] (c)--(x)--(d) (x)--(y)--(z)--(d);
  \foreach \v in {a,b,c,d} \node[retained] at (\v) {$\v$};
  \node[second] at (x) {$x$};
  \foreach \v in {y,z} \node[first] at (\v) {$\v$};
  \node[font=\footnotesize, text=coreteal] at (2.2,-1.3)
    {Remaining 3-core: $\{a,b,c,d\}$};
  \node[font=\footnotesize] at (2.2,-1.7)
    {Round 1: delete $y,z$.\quad Round 2: delete $x$.};
\end{tikzpicture}
\endgroup

  \caption{A fixed illustrative graph, not a random sample. Vertices $y,z$
    have degree two and disappear in the first simultaneous round. Vertex $x$
    initially has degree three, but loses $y$ and disappears in the second.
    The four mutually connected vertices $a,b,c,d$ form the final 3-core:
    its size is four and its fraction of the initial graph is $4/7$.}
  \label{fig:core-pruning}
\end{figure}

We use $M$ to distinguish the graph constructions defined below.
Their parameter $\lambda$ controls the mean number of memberships per vertex,
as specified in \cref{sec:triangle-model,sec:configuration-model}.
Let $G_{n,M}(\lambda)$ be the random input graph and define
\begin{equation}
  \begin{split}
    N_{n,M}(\lambda)&=|V(C_3(G_{n,M}(\lambda)))|,\\
    R_{n,M}(\lambda)&=N_{n,M}(\lambda)/n.
  \end{split}
  \label{eq:core-fraction}
\end{equation}
The denominator includes every initial vertex, including isolates.
When $n,M,\lambda$ are understood, we write simply $R$.
A core is empty if $R=0$. To separate substantial outcomes from small
nonempty cores, choose a fraction $\varepsilon_M\in(0,1]$ and set
\begin{equation}
  \begin{split}
    A_{n,M}(\lambda)&=\{R_{n,M}(\lambda)\geq\varepsilon_M\},\\
    p_{n,M}(\lambda)&=\Prob(A_{n,M}(\lambda)).
  \end{split}
  \label{eq:substantial-core}
\end{equation}
Thus survival means at least $\lceil n\varepsilon_M\rceil$ vertices remain;
equality with the cutoff is included. The baseline analyses use
$\varepsilon_M=0.05$ for all three graph conventions, with $0.02$ and $0.08$
as sensitivity checks in the triangle comparison. These are analysis choices
applied after pruning; they do not change the graph or its core.
At finite $n$, $0<R<\varepsilon_M$ means a nonempty core below the chosen
cutoff. A \emph{macroscopic} core instead describes growth with $n$:
its size remains a positive fraction of the input size. A single
finite-graph cutoff cannot establish that asymptotic property.

An unconditional mean or distribution includes all outcomes. The
conditional distribution keeps only substantial ones:
\begin{equation}
  F^+_{n,M}(r;\lambda)
  =\Prob\!\left(R_{n,M}(\lambda)\leq r\mid A_{n,M}(\lambda)\right),
  \label{eq:conditional-core-law}
\end{equation}
for $r\in[0,1]$, provided $p_{n,M}(\lambda)>0$.
Its mean and standard deviation describe sizes among survivors, whereas
$p_{n,M}$ describes how often they occur. The unconditional mean also
includes below-cutoff cores, which need not be empty. In an empirical
sample, estimate $p_{n,M}$ by the proportion meeting the cutoff and $F^+_{n,M}$
by the proportion of those survivors whose fractions are at most $r$.
With no survivors, conditional estimates are unavailable, not zero;
a sample standard deviation requires at least two. Existing baseline
size summaries are unconditional; conditional distributions are the separate
target of the finite-size investigation.

\subsection{Poisson triangle communities}
\label{sec:triangle-model}

A triangle community connects each of three individuals to the other two.
Combining such communities allows a vertex to belong to several triangles.
It also connects neighbours to one another, introducing the local clustering
discussed in the introduction. The construction is a triangle specialization
of random intersection graphs with communities~\cite{hofstad2021communities}.
The following specifies the finite simple graph used in the experiments.

Fix $n\geq1$ individuals, labelled $1,\ldots,n$, and a dimensionless parameter
$\lambda\geq0$. First propose independent membership counts
$\widetilde D_1,\ldots,\widetilde D_n$ with
\begin{equation}
  \Prob(\widetilde D_i=r)
  =e^{-\lambda}\frac{\lambda^r}{r!},
  \label{eq:poisson-memberships}
\end{equation}
for integers $r\geq0$. Here $\lambda$ is the mean number of membership places per individual
\emph{before conditioning}. For $\lambda=0$, every count is zero.
If the proposed total is not divisible by three, discard the entire vector
and draw again. Denote the accepted counts by $D_1,\ldots,D_n$ and their
total by $L=\sum_iD_i$. They have the joint Poisson law conditioned on
$L\equiv0\pmod3$. They are therefore not exactly independent Poisson variables
at finite $n$, and their marginal mean need not be exactly $\lambda$.
No single count is corrected and no leftover memberships are discarded.
The implementation aborts with an error after $100\,000$ unsuccessful
proposals; it never substitutes a different construction.

There are now $m=L/3$ communities, each with three labelled places.
To assign individuals to them:
\begin{enumerate}
  \item Make $D_i$ distinct membership tokens labelled $i$ for each
    individual $i$. Such a token, often called a membership stub,
    represents one place still to be assigned.
  \item Uniformly shuffle the $L$ tokens. Fill the three places of community
    one with the first three tokens, community two with the next three,
    and so on. Thus every matching of labelled tokens to labelled places
    has the same probability.
  \item In each triple, join every pair of distinct individuals, adding an
    edge only if that pair is not already connected.
\end{enumerate}
If $L=0$, there are no communities or edges, but all $n$ vertices remain.
Uniformity applies to token assignments, not to the resulting simple graphs.

The grouping does not require three different individuals. For example,
$(i,i,j)$ creates only the edge $\{i,j\}$, while $(i,i,i)$ creates none.
Repeated triples are also allowed. A connection from a vertex to itself
is omitted, and every repeated edge is merged into one, both within and
between communities. This conversion from membership places to edges is the
simple-graph \emph{projection}; the unmerged membership records are retained.
\Cref{fig:triangle-construction} shows both overlap and repeated places.
Equivalently, one can draw an incidence graph with one type of vertex for
individuals and another for communities, joining them once per assigned
place. Repeated memberships give repeated incidence edges. Core pruning is
performed on the projected graph of individuals, not on this incidence
graph, and does not require a prescribed number of whole triangles to survive.

\begin{figure}[b]
  \centering
  % Three prescribed groups of slots illustrate projection onto a simple graph.
% Two groups yield overlapping triangles. The repeated a in the third group
% contributes neither a loop nor a second copy of the already present edge ab.
\begingroup
\definecolor{motifteal}{RGB}{33,112,117}
\definecolor{slotgray}{RGB}{105,113,118}
\definecolor{slotorange}{RGB}{182,112,49}
\begin{tikzpicture}[
  x=1.15cm, y=1cm,
  vertex/.style={circle, minimum size=5.5mm, inner sep=0pt, font=\small,
    draw=motifteal, fill=white, text=motifteal, line width=0.8pt},
  slot group/.style={rectangle, rounded corners=2pt, minimum width=1.35cm,
    minimum height=5.5mm, inner sep=3pt, font=\small,
    draw=slotgray!70, fill=slotgray!5},
  edge/.style={draw=motifteal, line width=0.85pt}
]
  \node[font=\footnotesize] at (0,1.55) {Slot groups};
  \node[slot group] at (0,1) {$(a,b,c)$};
  \node[slot group] at (0,0.25) {$(c,d,e)$};
  \node[slot group, draw=slotorange, fill=slotorange!8] at (0,-0.5)
    {$(a,a,b)$};
  \draw[->, draw=slotgray, line width=0.8pt] (0.8,0.3)--(1.6,0.3);
  \node[font=\footnotesize] at (3.3,1.55) {Simple graph};
  % The resulting edge set is {ab,ac,bc,cd,ce,de}; c is shared.
  \coordinate (a) at (2.1,1);
  \coordinate (b) at (2.1,-0.5);
  \coordinate (c) at (3.4,0.25);
  \coordinate (d) at (4.7,1);
  \coordinate (e) at (4.7,-0.5);
  \fill[motifteal!7] (a)--(b)--(c)--cycle;
  \fill[motifteal!7] (c)--(d)--(e)--cycle;
  \draw[edge] (a)--(b)--(c)--cycle (c)--(d)--(e)--cycle;
  \foreach \v in {a,b,c,d,e} \node[vertex] at (\v) {$\v$};
  \node[font=\footnotesize, text=slotorange, align=center] at (2.05,-1.2)
    {The repeated $a$ adds no loop;\\$ab$ is already present.};
\end{tikzpicture}
\endgroup

  \caption{A fixed membership assignment and its simple projection.
    The triples $(a,b,c)$ and $(c,d,e)$ create two triangles sharing $c$.
    The triple $(a,a,b)$ creates no loop and adds no extra copy of $ab$.
    Membership counts are $(D_a,D_b,D_c,D_d,D_e)=(3,2,2,1,1)$, whereas
    distinct-neighbour degrees are $(2,2,4,2,2)$. This is an illustration
    of the construction, not a simulated realization.}
  \label{fig:triangle-construction}
\end{figure}

Each membership offers two edge ends, so write $K_i=2D_i$ for the
\emph{nominal} degree, before loops and repeated edges are removed.
If $\kappa_i$ is the actual number of distinct neighbours after projection,
then exactly
\begin{equation}
  \kappa_i\leq\min\{n-1,2D_i\}.
  \label{eq:triangle-degree-bound}
\end{equation}
Even a vertex with many memberships can have fewer neighbours because
some places repeat the same individual. Actual finite degrees can be odd.

For fixed $\lambda>0$ as $n\to\infty$, the divisibility condition has a
vanishing effect on any fixed number of membership counts.
Coincidences among the neighbours of a uniformly chosen vertex also become
unlikely. Its degree therefore converges in distribution to
\begin{equation}
  \begin{gathered}
    K=2D,\qquad D\sim\operatorname{Poisson}(\lambda),\\
    \pi_{2r}:=\Prob(K=2r)=e^{-\lambda}\lambda^r/r!,\\
    \pi_{2r+1}=0,\qquad r=0,1,\ldots.
  \end{gathered}
  \label{eq:limiting-degree-law}
\end{equation}
This is the triangle case of the community degree law in
\cite[Corollary~2.9 and the erasure discussion in Section~2.4]{hofstad2021communities}.
It gives $\E[K]=2\lambda$ and $\Var(K)=4\lambda$.
In particular, it is twice a Poisson variable, rather than a Poisson variable
of mean $2\lambda$. These are limiting degree statements, not identities
for the finite projected degrees or a theorem about the final 3-core.

\subsection{The degree-matched configuration model}
\label{sec:configuration-model}

The triangle construction joins neighbours in prescribed groups. To compare
it with connections formed without those groups, give each vertex a number
of individual edge ends and randomly join the ends in pairs. This is the
configuration model (CM)~\cite[Section~2]{janson2007kcore}.
For the implemented benchmark, independently draw
$D^{\mathrm{CM}}_1,\ldots,D^{\mathrm{CM}}_n$ from
$\operatorname{Poisson}(\lambda)$ and prescribe degrees
$K^{\mathrm{CM}}_i=2D^{\mathrm{CM}}_i$.
These counts are unconditioned: the total degree is automatically even.
Thus each prescribed degree has exactly the law in
\cref{eq:limiting-degree-law}, at every $n$.

Create $K^{\mathrm{CM}}_i$ distinct edge ends, or half-edge stubs, at vertex $i$.
Uniformly shuffle all ends and pair consecutive entries. A pair labelled
$(i,j)$ becomes an undirected edge, and every end is used.
Every matching of the labelled ends is equally likely; the resulting graphs
need not be equally likely. A general prescribed sequence with odd total
degree is rejected with an error, rather than repaired. That case cannot
arise for these twice-Poisson degrees.

Two ends of the same vertex form a \emph{self-loop}. Different pairs can
join the same two vertices, forming \emph{parallel edges}. The retained
CM graph keeps all of these, so it is a multigraph. Its degree counts
incident ends: a loop contributes two and each parallel edge contributes
separately. The degree at vertex $i$ is therefore exactly $K^{\mathrm{CM}}_i$.
For core pruning, keep a list of neighbours at each vertex, repeating a
neighbour once per edge to it and listing the vertex itself twice per loop.
The length of this adjacency list is its degree. On removing a vertex, mark it absent
first, then reduce each remaining neighbour's degree once per edge to it.
Its loops affect no other vertex. Queue a neighbour when its degree
crosses from three to two. The definition and order-independence argument
in \cref{sec:core-and-observables} apply with this degree convention.
It need not mean three distinct neighbours: two vertices joined by three
parallel edges already form a multigraph 3-core.

From the \emph{same} pairing we also obtain an erased simple CM graph:
delete every loop and retain one edge per unordered pair of distinct vertices.
Prune this graph separately using distinct-neighbour degrees.
Erasing can reduce degrees and core size, and can produce odd degrees.
It is different from rejecting an entire pairing until a simple graph is
obtained; in particular, its degrees need not equal the prescribed sequence.
The retained multigraph is the direct benchmark for the CM theory in
\cref{sec:configuration-theory}. The erased version checks the finite effect
of imposing the simple-graph convention used for triangles.

In the existing comparison, triangle and CM graphs use the same $n$ and
$\lambda$, with independently generated membership vectors and matchings.
Only the retained and erased CM outcomes share a pairing.
Here \emph{degree-matched} means a common limiting degree distribution,
not equal finite degree sequences, empirical histograms, or mean degrees.
A planned additional control would use one triangle-conditioned vector $D$
for both constructions and give the CM degrees $2D_i$.
It would match nominal degrees exactly, but still not generally match the
triangle graph's projected degrees. This shared-vector control is not part
of the existing baseline, and its conditioning differs from the unconditioned
CM sampler above.

To record finite departures from the common limiting law, let $\kappa_i$ be
the degree of vertex $i$ in the graph being measured: incident ends for the
retained CM and distinct neighbours for either simple graph.
The broad comparison computes all these degrees, including zeros.
It saves their histogram, mean, and population variance:
\begin{equation}
  \begin{split}
    H(j)&=\#\{i:\kappa_i=j\},\qquad j=0,1,\ldots,\\
    \bar\kappa&=\frac1n\sum_{i=1}^n\kappa_i,\qquad
    v_\kappa=\frac1n\sum_{i=1}^n(\kappa_i-\bar\kappa)^2.
  \end{split}
  \label{eq:degree-diagnostics}
\end{equation}
It also records the fraction $\sum_{j\text{ odd}}H(j)/n$ and the nominal
mean $n^{-1}\sum_iK_i$, using the corresponding model's nominal degrees.
For either construction, let $\ell$ count loops in the edge list before
erasure and $b$ count non-loop edge copies beyond the first per vertex pair.
For triangles, this list contains the three pairs of places from each
community. Erasure loses exactly $\ell+b$ edges, reducing total degree by
$2(\ell+b)$. These diagnostics quantify projection loss.

Clustering is measured on the simple graph in all cases. If it contains
$\tau$ triangles and has distinct-neighbour degrees $s_i$, define
\begin{equation}
  \mathcal T=\frac{3\tau}{\sum_{i=1}^n\binom{s_i}{2}},
  \label{eq:simple-transitivity}
\end{equation}
with $\mathcal T=0$ when the denominator is zero.
The denominator counts pairs of neighbours about a common vertex.
Each triangle closes three such pairs, so $\mathcal T$ is their connected
fraction, called transitivity. Both CM versions share this diagnostic because
they share their simple projection; it is measured before core pruning.

This benchmark preserves the limiting number of connections per vertex
while removing the requirement that they arrive in triangle communities.
The CM can still contain triangles formed by chance. Degree matching does
not control every correlation between vertices or every dependence created
by communities. Differences between the models therefore reflect those
structural changes together, rather than identifying an effect of clustering
alone. The next section develops the respective predictions for the measured
core fractions, keeping their different theoretical justifications explicit.
